\documentclass[11pt]{article}

\usepackage[margin=1in]{geometry}
\usepackage{amsmath,amssymb,amsthm}
\usepackage{graphicx}
\usepackage[hidelinks]{hyperref}

\newtheorem{theorem}{Theorem}
\newtheorem{remark}{Remark}

\newcommand{\norm}[1]{\left\lVert #1 \right\rVert}
\newcommand{\R}{\mathbb{R}}

\title{A Scalar Counterexample to a Uniform Property \texorpdfstring{$(S)$}{(S)} Modulus Question}
\author{Automated Math Discovery Smoke Test}
\date{June 8, 2026}

\begin{document}
\maketitle

\section*{Source and Open Question}

This packet concerns Open Problem (3) from:

\begin{quote}
William B. Johnson and Tomasz Kania, \emph{Uniform Property (S)},
arXiv:2602.09106.
\end{quote}

The source location is Section 5, Open problems, page 19 of the source
PDF; in the parsed source this is line 1071 of \texttt{main.tex}.

\begin{center}
\includegraphics[width=\textwidth]{figures/open_problem_crop.png}
\end{center}

Problem (3) asks:

\begin{quote}
Is $U(d,a)$ always smaller than $da$ (up to a constant)?
\end{quote}

We give a negative answer to the statement as written, in the real
Banach-space category, even if one restricts to spaces with uniform
property $(S)$.

\section*{Definitions}

For a normed space $X$, Johnson and Kania define
\[
  U_X(x,y;a)
  =
  \sup_{\norm{z}\le a}
  \left|\,\norm{x+z}-\norm{y+z}\,\right|
\]
for $x,y\in S_X$ and $a>0$, and
\[
  U_X(d;a)
  =
  \inf_{\substack{x,y\in S_X\\ \norm{x-y}\ge d}}
  U_X(x,y;a).
\]
The open question asks whether there is a universal constant $C$ such
that
\[
  U_X(d;a)\le Cda
\]
for all relevant spaces $X$, all $0<d<2$, and all $a>0$.

\section*{Counterexample}

\begin{theorem}
There is no universal constant $C$ such that
\[
  U_X(d;a)\le Cda
\]
for every real Banach space $X$ with uniform property $(S)$, every
$0<d<2$, and every $a>0$.
\end{theorem}

\begin{proof}
Let $X=\R$ with its usual norm. Then
\[
  S_X=\{-1,1\}.
\]
Consequently, if $0<d<2$, the only admissible ordered pairs
$(x,y)\in S_X\times S_X$ with $\norm{x-y}\ge d$ are $(1,-1)$ and
$(-1,1)$. Therefore
\[
  U_{\R}(d;a)=U_{\R}(1,-1;a).
\]

We compute this quantity. For $z\in\R$,
\[
  \left|\,|1+z|-|-1+z|\,\right|\le 2
\]
by the reverse triangle inequality. If $0<a\le 1$, then for
$z\in[-a,a]$ one has $1+z\ge 0$ and $1-z\ge 0$, so
\[
  |1+z|-|-1+z|
  =
  (1+z)-(1-z)
  =
  2z.
\]
Taking absolute values and then the supremum over $|z|\le a$ gives
\[
  U_{\R}(d;a)=2a
  \qquad (0<d<2,\; 0<a\le 1).
\]
In fact the full scalar formula is
\[
  U_{\R}(d;a)=\min\{2a,2\},
  \qquad 0<d<2,\ a>0.
\]
In particular, $U_{\R}(d;a)>0$ for every $0<d<2$ and every $a>0$, so
$\R$ has uniform property $(S)$.

Now suppose that a universal constant $C$ existed. Applying the proposed
bound to $X=\R$ and $0<a\le 1$ would give
\[
  2a=U_{\R}(d;a)\le Cda
\]
for every $0<d<2$. Since $a>0$, this implies
\[
  2\le Cd
\]
for every $0<d<2$, which is impossible as $d$ tends to zero.
\end{proof}

\section*{Verification Notes}

The proof is exact. The script
\texttt{code/check\_scalar\_modulus.py} numerically compares the formula
\[
  U_{\R}(d;a)=\min\{2a,2\}
\]
against a fine grid maximization of
\[
  \left|\,|1+z|-|-1+z|\,\right|
  \qquad (|z|\le a)
\]
for several values of $a$. This is only a sanity check; the scalar
calculation above is the proof.

\section*{Limitations and Human Review Guidance}

This is a literal counterexample to Problem (3) as printed. It is also a
small edge case: the obstruction is that the real one-dimensional unit
sphere has no pairs at distance $d$ for $0<d<2$ except the antipodal
pair, whose distance is $2$. Thus the infimum defining $U(d;a)$ cannot
probe genuinely small-distance geometry.

The result should therefore be read as identifying a missing hypothesis
in the printed question. A more substantive modified problem might ask
for the same upper bound in infinite-dimensional spaces, in spaces with
connected unit sphere, or under a condition guaranteeing admissible
pairs at every small distance. This packet does not resolve those
strengthened versions.

\section*{Novelty Check}

No external literature search beyond the source paper and the local
pipeline corpus was performed. The argument is elementary enough that it
may be known informally; the claim here is only that it answers the
printed statement exactly.

\end{document}
